根据实验要求，目标密文序号由小组成员学号最后一位之和对 10 取模得到。本小组抽到的目标密文序号为 5。实验材料中给出的 5 号密文属于 Key\_2 加密所得密文，因此本实验在攻击过程中使用 \verb|dec_2| 解密接口进行 Padding Oracle 查询。

\subsection{目标密文}

实验材料中给出的 5 号目标密文为：

\begin{lstlisting}[breaklines=true, basicstyle=\small\ttfamily]
46307250616464316e674f7261636c339ae0735429869542efc40dcdc3f4c170649463f719c5ddf4ce8c6d1ef0e5d41ac5d137629e3fe1340cfaad7e21d65d14
\end{lstlisting}

该密文采用 AES-128-CBC 模式加密，并使用 PKCS\#7 填充。由于 AES 的分组长度为 16 字节，因此需要按照每 16 字节，即每 32 个十六进制字符，对密文进行分组。

\subsection{密文分组}

目标密文的整体结构为：

\begin{equation}
    IV \parallel C_1 \parallel C_2 \parallel C_3
\end{equation}

其中，$IV$ 为初始向量，$C_1,C_2,C_3$ 为三个密文分组。按照 16 字节分组后得到：

\begin{table}[H]
    \centering
    \caption{5号目标密文分组结果}
    \begin{tabular}{|c|l|}
        \hline
        分组 & 十六进制表示 \\
        \hline
        $IV$ & \texttt{46307250616464316e674f7261636c33} \\
        \hline
        $C_1$ & \texttt{9ae0735429869542efc40dcdc3f4c170} \\
        \hline
        $C_2$ & \texttt{649463f719c5ddf4ce8c6d1ef0e5d41a} \\
        \hline
        $C_3$ & \texttt{c5d137629e3fe1340cfaad7e21d65d14} \\
        \hline
    \end{tabular}
\end{table}

程序运行时首先对目标密文进行分组，输出结果如下：

\begin{lstlisting}[breaklines=true, basicstyle=\small\ttfamily]
IV = 46307250616464316e674f7261636c33
C1 = 9ae0735429869542efc40dcdc3f4c170
C2 = 649463f719c5ddf4ce8c6d1ef0e5d41a
C3 = c5d137629e3fe1340cfaad7e21d65d14
\end{lstlisting}

在 CBC 解密过程中，$IV$ 可以视为第 0 个密文分组，即 $C_0 = IV$。因此三个明文分组的恢复关系分别为：

\begin{align}
    P_1 &= D_K(C_1) \oplus IV, \\
    P_2 &= D_K(C_2) \oplus C_1, \\
    P_3 &= D_K(C_3) \oplus C_2.
\end{align}

后续攻击过程的目标就是分别恢复 $D_K(C_1)$、$D_K(C_2)$ 和 $D_K(C_3)$ 的中间状态，再与对应的前一分组异或，最终得到完整明文。